\SourceLeaf{369}%
wenn auch nur für eine kugelförmige Erde, und wie man den Abstand einzelner Punkte auf der Erdoberfläche voneinander bestimmt. Überhaupt hätten wir gewünscht, dass die sphärische Trigonometrie nicht so vollständig ausgeklammert worden wäre. Auch bei der sphäroidischen Gestalt der Erde hätte wohl \textit{deutlicher} herausgearbeitet werden können, wie der Begriff der geografischen Breite gegenüber dem Fall der Kugel geändert werden muss und wie die relative Lage zum Erdäquator, zur Erdachse und zum Erdmittelpunkt von dieser Breite abhängt. Doch das sind Kleinigkeiten, die dem Gesamtwert des Buches keinen Abbruch tun. Der Verfasser kann sie leicht berücksichtigen, falls eine neue Auflage nötig werden sollte, wie es bei einem für den Unterricht sehr empfehlenswerten Buch durchaus möglich ist.

